\section{直接股票排除与启用边界}

本案当前策略不持有直接股票，也不持有期货。直接股票初始权重为0\%，不设发行人名单、目标价、历史收益率或现金分配假设；期货初始权重同为0\%，不作为替代配置工具。

\begin{exhibitbox}[当前边界：直接股票与期货均不进入初始策略]
\small
\begin{tabularx}{\exhibitboxwidth}{L{3.0cm}P P}
\toprule
\textbf{风险资产类别} & \textbf{初始权重} & \textbf{当前处理} \\
\midrule
直接股票 & 0\% & 不建立持仓、不形成标的排序，也不纳入现金、估值或执行结论。 \\
期货 & 0\% & 不使用杠杆或衍生品替代现货基金配置。 \\
\bottomrule
\end{tabularx}
\end{exhibitbox}

\section{排除的原因}

现行材料没有形成可独立复核的发行人级估值、现金流、税务、集中度和流动性证据包。因此，任何直接股票敞口都会把未完成的发行人判断带入组合；在这些判断完成前，零权重是策略约束，而不是等待选择的观察池。

\section{未来版本的启用门槛}

以下为未来经独立复核的策略修订所必须满足的边界，不构成当前启用授权：

\begin{enumerate}
  \item 对每一发行人完成独立、时点明确的估值审查，并说明估值结论的证据边界；
  \item 核验现金流、资本开支、分配覆盖和资产负债表约束，不以历史现金分配代替前瞻判断；
  \item 按适用账户和持有期完成税务审查，不将基金规则直接套用于直接股票；
  \item 穿透检验与现有权益基金的持仓重叠、行业／集团集中度及相关性；
  \item 完成可交易性、容量、停复牌、成本和退出路径审查，并在实施日重新核验；
  \item 取得独立复核和正式策略修订记录后，方可讨论是否启用。
\end{enumerate}

即使未来所有门槛通过，直接股票总额仍不得超过组合的15\%，单一发行人不得超过5\%。任何拟议敞口须受同等或更严格的集中度、流动性和治理约束；未通过、过期或缺失的任一门槛均意味着直接股票权重保持0\%。

\sourcenote{本章仅陈述当前策略边界。已退役的历史行情、分配和期货资料均不是本章输入，且不得支持目标、配置、估值、现金或执行结论。}
